To estimate the effect that wartime relocation and internment had on wages, I need to compare the outcomes of likely internees before and after 1942 against a reasonable control group.
The vast number of Japanese Americans in the US at the time were interned, so the size of the potential control group of Japanese living outside of the West Coast is quite small.
For this reason, I follow \citeauthor{arellano-boverDisplacementDiversityMobility2022} by including Chinese Americans in my main sample.

Chinese immigrants entered the US West Coast at a similar period as the major wave of Japanese immigration.
Japanese workers were also seen as a substitute for the cheap Chinese labor for work on the transcontinental railroads and in agriculture 
\citep{chenImpactSkillBasedImmigration2015}.
Chinese Americans were not included in any of the relocation or internment orders following Executive Order 9066,
but in the absence of internment, their wages and employment would likely have continued to evolve similarly to the West Coast Japanese population.
For these reasons, Chinese make a natural comparison group to recreate the counterfactual earnings outcomes interned Japanese would have had if they had never been interned.

Although all persons of Japanese descent living on the West Coast were ordered to evacuate outside of the Exclusion Area
(and no Chinese persons ordered to evacuate),
the ability of some Japanese to self-evacuate means some were able to avoid being sent to any of the ten internment camps.
The \cite{dewittjohnl.FinalReportJapanese1943} puts the number of voluntary evacuees at 4,889.
Other studies which treat all Japanese living on the West Coast in 1940 as the treatment group
ignore this number as well as any voluntary emigration by Japanese out of the West Coast to other states.
In order to estimate a more accurate probability of the actual probability that any individual Japanese census respondent was a true internee,
I create a measure of predicted internee status.

\subsection{Predicted Internment Status}\label{sec:data-intern}

The WRA records give a complete list of all individuals who were interned during World War II.

To determine the probability that someone who appears in the 1940 census would go on to be interned in 1942, I use the method of \cite{arellano-boverDisplacementDiversityMobility2022}.
Per his application of Bayes' Rule, the probability that an individual $i$ in the demographic group $z$ who lived in county $c$ in 1942 was interned ($Pr(I_{i}=1|z_i,c_{i})$
could be determined from the proportion of the whole population who were interned ($Pr(I_{i}=1)$,
the proportion of the population with demographics $z$ who lived in county $c$ in 1942 ($Pr(z_i,c_{i}^{1942}$),
and the probability of a known internee having the same demographics and county in 1942 ($Pr(z_i,c_i^{42}|I_i=1)$).

$$Pr(I_i=1|z_i,c_i^{42}) = \frac{Pr(z_i,c_i^{42}|I_i=1) \cdot Pr(I_i=1)}{Pr(z_i,c_i^{42})}$$

However, because my target sample of census respondents were only asked about their county of residence in the year of the census and not in 1942, $c_i^{1940}$ is used as a proxy for census respondents' 1942 locations.

The estimator for the true internment status of a census respondent is:

\begin{equation}\label{eq:predint}
\hat{E}[I]_{zc} =
\frac{n(i | z_i=z, ~ c_{i}^{42}=c)}{n(j | z_j=z, ~ c_{j}^{40}=c)}
\end{equation}

This will produce a biased estimate of true internment status if there were unobserved changes in population in each county between 1940 and 1942. For example, for a county that experienced a net outflow of residents, then the ratio of residents in 1940 vs 1942 would be greater than 1 and we would have an estimated internment probability that overstates the likelihood of anyone who lived in that county in 1940 would be interned.

One source of differences between the number of 1940 residents and the number of internees who came from each county is that the Army initially allowed Japanese persons to choose to move out of the exclusion areas before they gathered the remaining people into assembly centers.

\subsection{Internment's Effect on Earnings}

To measure the effect of internment on earnings, my main analysis follows the following specification.

\begin{equation}\label{eq:mainreg}
  Y_i = \alpha_i + \beta Pr(I)_i + \mathbb{X}_i \Gamma + \epsilon_i
\end{equation}

The main outcome of interest which I use for $Y_i$ is the reported dollar earnings of individual $i$ in the 1950 census.
$\alpha_i$ represent the state fixed effects for the state in which individual $i$ is living in 1950.
The vector of individual-level controls in $\mathbb{X}$ include age, sex, and individual $i$'s earnings from the 1940 census,
adjusted to 1950 dollars.

The probability of individual $i$ being an internee, $Pr(I)_i$,
is replaced by my estimated version from equation \ref{eq:predint}
using individual $i$'s reported race in 1940 as the only demographic $z$.
The main causal effect of interest is represented by the estimator $\beta$,
which represents the number of dollars of income which an individual is estimated to have gained (or lost)
when going from a non-internee to a likely internee.